% Источники: exam/materials/моделирование.pdf, раздел 23; https://github.com/HumsterProgrammer/iu9-notes/blob/main/8sem/Моделирование/Моделирование_лекция-04_19-02-26.md.
\section{Модельное время. Подходы к дискретизации времени.}

\subsection*{Виды времени}

В имитационном моделировании различают три вида времени:
\begin{itemize}
  \item \textbf{Реальное время} — астрономическое время функционирования реального объекта.
  \item \textbf{Вычислительное время} — время, затраченное ЭВМ на расчёт модели (машинное время).
  \item \textbf{Модельное время} $\tau$ — внутренние «часы» модели; шкала времени, в единицах которой описывается поведение объекта.
\end{itemize}

Масштаб модельного времени относительно реального может быть произвольным: моделирование быстрее, медленнее или в реальном масштабе времени.

\subsection*{Подходы к дискретизации модельного времени}

\textbf{1. Постоянный шаг ($\Delta\tau = \mathrm{const}$).}

Модельное время продвигается равномерными шагами $\Delta\tau$: $\tau_{k+1} = \tau_k + \Delta\tau$.
\begin{itemize}
  \item Просто реализуется; используется при подходе на основе активностей.
  \item Недостаток: при редких событиях выполняется много «холостых» шагов; при частых — возможен пропуск событий между шагами.
\end{itemize}

\textbf{2. Переменный шаг.}

Модельное время скачком переходит к моменту следующего ближайшего события:
\[
  \tau_{k+1} = \min_{i}\bigl\{\tau^i_{\text{след}}\bigr\},
\]
где $\tau^i_{\text{след}}$ — запланированное время очередного события $i$-го компонента.
При выборе модельного времени с переменным шагом шаг выбирается по моменту следующего изменения состояния. Если два и более событий для любой компоненты попадают в один квант модельного времени, соответствующий промежуток делится пополам; при необходимости деление повторяется.

В источнике отмечено, что подход с переменным шагом имеет слабо формализованную процедуру перехода от реального времени к модельному времени, что приводит к специальным способам организации имитации.

\clearpage
